\scanpage{14}
\handwritten{\mspage{5}}
% archive sheet number: 13
% The top typed insertion is attached by a pen guide after Onnsuppose.
Onnsuppose \inserted{a) et b) vérifiés pour le morphisme $i_{\bar s}:X_{\bar s}\handwritten{\longrightarrow}X_{\bar s}$ induit} \deleted{que le morphisme induit} pour lesfibres en $\bar s$ (point géométrique \deleted{fixé} de $S$\deleted{, qu'on a intérêt à prendre} au dessus da point \uline{générique} de $S$). \deleted{si $S$ est irréductible pour avoir le maximum de chances de \unsure{pouvoir que|pourvoir que} l'hypothèse sérieuse b). Soit $z$ un cycle a} \inserted{Soit} $z\handwritten{\in}Z^n(Y)$ tel que le cycle induit sur $Y_{\bar s}$ soit algébriquement équivalent à zéro, et à fortiori cohomologiquement équivalent à zéro. Donc da classe
\[
u(z)\handwritten{\in}H^1(S,R^{2n-1}f_{\handwritten{*}}(\underline{Z}_{\unsure{\chi|\ell}}(n))
\]
% The displayed H-group has no third closing parenthesis in the source.
est définie. Ceci posé, je dis que \uline{cette classe est \deleted{une classe de torsion} contenue dans} $i^{\handwritten{*}}(H^{\deleted{\unsure{0}}1}(S,R^{2n-1}\deleted{\unsure{f|g}}\handwritten{g}_{\handwritten{*}}(\unsure{\underline{\deleted{Z}\handwritten{Q}}|\underline{Z}}_{\chi}(n)))$ \uline{modulo une classe de torsion}, du moins si on \deleted{suppose $S$ japonais de dimension 1, ou même} \deleted{\inserted{On applique 2) à $i_{\bar s}$}} \deleted{\inserted{si,}} se permet de remplacer $S$ par un ouvert non vide(restriction \deleted{sans doute}\inserted{probablement} inutile, mais cet énoncé nous suffira pour ce qu'on veut en faire).
% The inline i*(H-group) likewise lacks its outermost closing parenthesis.
% The exact attachment of the already-cancelled On applique interline is not recoverable.

Pour le voir, on applique 2) au morphisme $i_{\bar s}$. La courbe $C$ et le cycle $Z''$ peuvent \unsure{d|\deleted{d}s|\deleted{s}d}e défini\unsure{t|\deleted{t}r|\deleted{r}t} déjà sur une extension finie $K$ de $k(s)$, Un argument de trace permet de supposer que $K=k(s)$, et \deleted{quitte à prendre l'adhérence du cycle} on peut supposer que $\unsure{C|Z|\mathbb{C}}$ \deleted{et $Z$} est une courbe propre et lisse relative sur $S$ lui-même, et $Z''$ un cycle sur $\unsure{X|\mathbb{X}|Y}\mathbin{\mathrm{x}}_S C$, \deleted{Alors} et que la formule $z=i^{\handwritten{*}}(z')+z''$ est valable sur $X\mathbin{\mathrm{x}}_S C$. Elle implique une décomposition analogue de $u(z)$, qui est la somme de $i^{\handwritten{*}}(u(z'))$ et de $u(z'')$, de sorte qu'il suffit de prouver que $u(z'')$ est nul. Or $u(z'')$ induit par \deleted{d}construction 2) l'homomorphisme nul $H^1(C_{\bar s})\handwritten{\longrightarrow}H^{2n-1}(Y_{\bar s})$. Im résulte donc bien de 1) que $u(z'')$ est nul, cqfd.
% Both product capitals were individually inspected; the second repeats the first candidate.
% The C/Z overprint has no securely recoverable chronological order.

4.\quad Soient $X$ un schéma projectif et lisse de dimension $m$\inserted{\handwritten{sur $k$ Alg. clos,}}, prenons un pinceau de Lefschetz, de façon à obtenir une variété éclatée non moins projective et lisse $\overset{\handwritten{\sim}}{X}$ et un morphisme
\[
\overset{\handwritten{\sim}}{f}:\overset{\handwritten{\sim}}{X}\handwritten{\longrightarrow}\deleted{S\mathrm{x}}P^1_k
\]
% The x after the cancelled S records one additional visible cancellation cell,
% not a recovered old lexical letter. The pen comma after clos and the typed comma
% after m are both visible; the insertion is placed before the latter.
% Unidentified narrow pen stroke touching the top of b in fibres (whole box
% [414,1899,433,1948]); its linguistic function is uncertain. The following figure
% reproduces that mark only, separately to avoid changing the word fibres.
\figure{\begin{tikzpicture}[x=1cm,y=1cm]
\draw[line width=.25pt] (.09,.20) -- (.11,.40) -- (.13,.62);
\draw[line width=.45pt] (.08,0) -- (.08,.12) -- (.06,.19);
\node[inner sep=0pt] at (.33,.31) {$\unsure{\phantom{x}}$};
\end{tikzpicture}}
dont les fibres sont les sections hyp\deleted{é}erplanes du pinceau. Soit $\unsure{S|\mathbb{S}|U}$ un ouvert non vide de $P^1_k$ au dessus duquel $\overset{\handwritten{\sim}}{X}$ soit lisse, et soit $\handwritten{\unsure{j|f}}:Y=\overset{\handwritten{\sim}}{X}\handwritten{|}S\handwritten{\longrightarrow\unsure{S}}$.
